\begin{theorem} \label{theorem:internal_and_external_direct_sums_of_vector_spaces_are_isomorphic}
    Let $F$ be a \CrefAndHyperrefIfExist{definition:field}{field}. Let $V$ be an \CrefAndHyperrefIfExist{definition:vector_space_over_a_field}{$F$-vector space}, and let $W_1, \dots, W_n$ be subspaces of $V$ whose internal direct sum $W = W_1 \oplus \cdots \oplus W_n$ exists.

    The map $\Phi: W_1 \oplus \cdots \oplus W_n \to W$ defined by
    $$\Phi(w_1, \dots, w_n) = w_1 + \cdots + w_n$$
    is an \CrefAndHyperrefIfExist{definition:morphism_of_vector_spaces}{isomorphism of $F$-vector spaces} between the external and internal direct sums.

\TextIfExists{definition:coproduct_of_modules_of_rings}{This is the vector-space special case of the canonical isomorphism between an internal direct sum of submodules and \Cref{definition:coproduct_of_modules_of_rings}.}
\TODO{AI generated, verify}
\end{theorem}
\begin{proof}
    \TODO{AI generated, verify}
    By \Cref{definition:direct_sum_of_vector_spaces_internal_and_external}, the external direct sum $W_1 \oplus \cdots \oplus W_n$ is the Cartesian product equipped with component-wise addition and scalar multiplication.

    We first show that $\Phi$ is a linear map. For $(w_1, \dots, w_n), (w'_1, \dots, w'_n) \in W_1 \oplus \cdots \oplus W_n$ and $c \in F$, we have
    $$\Phi(w_1 + w'_1, \dots, w_n + w'_n) = (w_1 + w'_1) + \cdots + (w_n + w'_n) = (w_1 + \cdots + w_n) + (w'_1 + \cdots + w'_n) = \Phi(w_1, \dots, w_n) + \Phi(w'_1, \dots, w'_n),$$
    and
    $$\Phi(cw_1, \dots, cw_n) = cw_1 + \cdots + cw_n = c(w_1 + \cdots + w_n) = c\Phi(w_1, \dots, w_n).$$
    Hence $\Phi$ is linear by \Cref{definition:morphism_of_vector_spaces}.

    We show that $\Phi$ is surjective. Let $w \in W$. By the definition of the internal direct sum (\Cref{definition:direct_sum_of_vector_spaces_internal_and_external}), every element of $W = W_1 + \cdots + W_n$ can be written as a sum $w_1 + \cdots + w_n$ with $w_i \in W_i$. Hence, $\Phi(w_1, \dots, w_n) = w$, so $\Phi$ is surjective.

    We show that $\Phi$ is injective by verifying it has trivial kernel. Suppose that $\Phi(w_1, \dots, w_n) = 0_V$. Then $w_1 + \cdots + w_n = 0_V$ with $w_i \in W_i$. By the uniqueness of the expression in an internal direct sum (\Cref{definition:direct_sum_of_vector_spaces_internal_and_external}), this implies that each $w_i = 0_{W_i}$. Hence, $(w_1, \dots, w_n) = (0_{W_1}, \dots, 0_{W_n})$, which is the zero vector of the external direct sum. Therefore, $\Phi$ has trivial kernel, so it is injective by \Cref{proposition:basic_facts_about_kernel_and_image_of_group_homomorphisms}\eqref{item:group_hom_is_injective_iff_kernel_is_trivial} (applying the group-theoretic statement to the underlying abelian groups of vector spaces).

    Hence, $\Phi$ is a bijective linear map, i.e., an isomorphism of $F$-vector spaces.
\end{proof}
